% !TEX program = xelatex
\documentclass[aspectratio=169,11pt]{beamer}

% ------------------------------------------------------------------
%  Persona-Driven Election Simulation  --  clean Beamer deck
% ------------------------------------------------------------------

% --- Block metropolis's automatic Fira Sans/Mono font check BEFORE the
% theme loads. Some MiKTeX/Windows setups report a false "font found"
% and then crash trying to build it (hbf2gf/miktex-maketfm errors).
% Forcing \IfFontExistsTF to always take the "not found" branch while
% metropolis loads sidesteps that entirely, regardless of local font
% cache state.
\usepackage{iftex}
\ifxetex
\usepackage{fontspec}
\else
\ifluatex
\usepackage{fontspec}
\fi
\fi
\makeatletter
\let\IfFontExistsTF@saved\IfFontExistsTF
\def\IfFontExistsTF#1#2#3{#3}% always take the "does not exist" branch
\usetheme{metropolis}                 % clean, minimal, no clutter
\let\IfFontExistsTF\IfFontExistsTF@saved
\makeatother

% Now explicitly set safe fonts that ship with every TeX install.
\ifxetex
\setsansfont{Latin Modern Sans}
\setmonofont{Latin Modern Mono}
\else
\ifluatex
\setsansfont{Latin Modern Sans}
\setmonofont{Latin Modern Mono}
\fi
\fi

\usepackage[utf8]{inputenc}
\usepackage[T1]{fontenc}
\usepackage{amsmath,amssymb}
\usepackage{booktabs}                 % nice tables
\usepackage{siunitx}                  % aligned numbers
\usepackage{graphicx}
\usepackage{tikz}
\usepackage{eurosym}

%: a couple of accent colors (subtle, not flashy) ---------------
\definecolor{accent}{HTML}{2A7DE1}
\definecolor{amber}{HTML}{C77A00}
\setbeamercolor{frametitle}{bg=black!88}
\setbeamercolor{progress bar}{fg=accent}

\metroset{block=fill}

% siunitx table setup
\sisetup{group-separator={,}, group-minimum-digits=4}

\newcommand{\mynote}[1]{\vspace{0.4em}\footnotesize\color{gray}#1}

%: placeholder box shown when an image file hasn't been added yet,
% so the deck still compiles while you're filling in real figures.
\newcommand{\missingimg}[2][0.7\textheight]{%
	\fbox{\parbox[c][#1][c]{0.92\linewidth}{\centering
			\color{gray}\itshape Missing image:\\ \texttt{\detokenize{#2}}}}}

%: helper: one big centered image, using as much space as possible
\newcommand{\bigimg}[1]{%
	\begin{center}
		\IfFileExists{#1}%
		{\includegraphics[width=0.98\textwidth,height=0.9\textheight,keepaspectratio]{#1}}%
		{\missingimg[0.7\textheight]{#1}}
\end{center}}

%: helper: two images side by side, left half / right half, no caption, using max space
\newcommand{\splitimg}[2]{%
	\begin{columns}[T]
		\column{0.5\textwidth}
		\centering
		\IfFileExists{#1}%
		{\includegraphics[width=0.98\linewidth,height=0.86\textheight,keepaspectratio]{#1}}%
		{\missingimg[0.7\textheight]{#1}}
		\column{0.5\textwidth}
		\centering
		\IfFileExists{#2}%
		{\includegraphics[width=0.98\linewidth,height=0.86\textheight,keepaspectratio]{#2}}%
		{\missingimg[0.7\textheight]{#2}}
\end{columns}}

\title{Persona Election Simulation}
\author{Jorge Andrés Pulgarín Barbosa}
\institute{University of Tübingen}
\date{}

\begin{document}
	
	% ================================================================
	\maketitle
	
	% ================================================================
	% SLIDE -- Kish allocation
	% ================================================================
	\begin{frame}{Kish Allocation: what does $K$ do?}
		K controls the trade-off between \textbf{equal} and
		\textbf{proportional} allocation:
		\[
		\boxed{\;n_h \;=\; n \cdot
			\frac{N_h^{\,K}}{\displaystyle\sum_j N_j^{\,K}}\;}
		\]
		\vspace{0.3em}
		\begin{columns}[t]
			\column{0.32\textwidth}
			\begin{block}{$K = 0$}
				$N_h^{0}=1$ \\
				Every region gets the \textbf{same} sample size.
			\end{block}
			\column{0.32\textwidth}
			\begin{block}{$0 < K < 1$}
				A \textbf{compromise} between equal and proportional.\\[0.3em]
				\small\color{amber}$K=0.67$ in Canadian survey design.
			\end{block}
			\column{0.32\textwidth}
			\begin{block}{$K = 1$}
				$n_h \propto N_h$ \\
				Fully \textbf{proportional} to population.
			\end{block}
		\end{columns}
	\end{frame}
	
	% ================================================================
	% SLIDE -- Canada 72,000 example (BIG)
	% ================================================================
	\begin{frame}{Canada: real 72{,}000-sample survey}
		A rescaling of the actual Canadian survey design with
		\textbf{\color{amber}$K=0.67$}:\quad
		$n_h = 72{,}000 \cdot N_h^{\,0.67} \big/ \sum_j N_j^{\,0.67}$
		\vspace{0.3em}
		\begin{center}
			{\renewcommand{\arraystretch}{1.18}\setlength{\tabcolsep}{12pt}
				\begin{tabular}{l r r r r}
					\toprule
					\textbf{Province} & \textbf{Population}
					& $K{=}0$ & \textbf{\color{amber}$K{=}0.67$} & $K{=}1$ \\
					\midrule
					Newfoundland \& Labrador & \num{433148}   & \num{7200} & \color{amber}\num{2378}  & \num{1108}  \\
					Prince Edward Island     & \num{118756}   & \num{7200} & \color{amber}\num{999}   & \num{304}   \\
					Nova Scotia              & \num{774025}   & \num{7200} & \color{amber}\num{3508}  & \num{1981}  \\
					British Columbia         & \num{3719183}  & \num{7200} & \color{amber}\num{10041} & \num{9517}  \\
					Ontario                  & \num{10083279} & \num{7200} & \color{amber}\num{17064} & \num{25797} \\
					Quebec                   & \num{7112070}  & \num{7200} & \color{amber}\num{13842} & \num{18196} \\
					\bottomrule
			\end{tabular}}
		\end{center}
		\mynote{\textit{Source: Access and Support to Education and Training Survey (ASETS), Statistics Canada.}}
	\end{frame}
	
	% ================================================================
	% SLIDE -- Our election allocation (BIG, no proportional column)
	% ================================================================
	\begin{frame}{From Kish to our election allocation}
		Kish's idea: balance competing objectives with a \textbf{single score per
			state}. We use \textbf{three drivers}, each turned into a share that sums
		to 1 across states:
		\vspace{0.6em}
		
		\begin{columns}[T]
			\column{0.5\textwidth}
			\begin{itemize}
				\item \textbf{Population} \\
				$W_{\text{Pop},d} = Pop_d \big/ \sum_h Pop_h$
				\item \textbf{Electoral votes} \\
				$W_{\text{EV},d} = EV_d \big/ \sum_h EV_h$
				\item \textbf{Competitiveness} \\
				$W_{\text{Comp},d}$  larger for states that were
				\emph{close} in past elections.
			\end{itemize}
			
			\column{0.5\textwidth}
			All three matter equally, so we combine them into one score:
			\[
			Z_d = \sqrt{\,W_{\text{Pop},d}^2 + W_{\text{EV},d}^2
				+ W_{\text{Comp},d}^2\,}
			\]
			\vspace{0.2em}
			
			Each state then receives a share of the \textbf{200{,}000} personas
			\textbf{proportional to $Z_d$}.
		\end{columns}
	\end{frame}
	
	% ================================================================
	% OUR EXAMPLE -- why the allocation changes by state (CA / PA / WY)
	% ================================================================
	\begin{frame}{Our example: why the allocation changes by state}
		Three states, three different drivers:
		\textbf{size}, \textbf{competition}, and \textbf{neither}:
		\vspace{0.3em}
		\begin{center}
			\small
			\begin{tabular}{l r r r r r r}
				\toprule
				\textbf{State} & \textbf{Pop.} & \textbf{Pop.\ w.}
				& \textbf{EV} & \textbf{EV w.} & \textbf{Comp.\ w.} & \textbf{$\sim$calls} \\
				\midrule
				California   & $\sim$39.5M & 0.1151 & 54 & 0.1004 & 0.0003 & $\sim$16{,}200 \\
				Pennsylvania & $\sim$13.0M & 0.0382 & 19 & 0.0353 & 0.0904 & $\sim$11{,}100 \\
				Wyoming      & \num{588753} & 0.0017 & 3 & 0.0056 & $\sim$0.0000 & 619 \\
				\bottomrule
			\end{tabular}
		\end{center}
		\vspace{0.2em}
		\[
		Z_{\text{CA}} = 0.1528 \qquad
		Z_{\text{PA}} = 0.1043 \qquad
		Z_{\text{WY}} = 0.0058 \qquad
		\textstyle\sum_d Z_d = 1.9917
		\]
		\begin{itemize}\small
			\item \textbf{California}: large population + many electoral votes;
			competitiveness contributes almost nothing.
			\item \textbf{Pennsylvania}: much smaller, but a historically tight
			swing state; competitiveness drives most of its allocation.
			\item \textbf{Wyoming}: small on all three, so it receives one of
			the smallest allocations.
		\end{itemize}
	\end{frame}
	
	% ================================================================
	% SLIDE -- Cost / scale
	% ================================================================
	\begin{frame}{200{,}000-person simulation: scale, time, cost}
		\begin{center}
			\begin{tabular}{c c c c}
				\textbf{200{,}000} & \textbf{140--150 min} & \textbf{144.6M} & \textbf{\color{accent}\$2.43--\$2.60} \\
				personas / calls & runtime & total tokens & our inference setup \\[0.2em]
				& & 114.2M in / 30.4M out & (\$7.38 at API pricing) \\
			\end{tabular}
		\end{center}
		\vspace{0.6em}
		
		\begin{center}
			{\renewcommand{\arraystretch}{1.25}\setlength{\tabcolsep}{14pt}
				\begin{tabular}{l r r r}
					\toprule
					\textbf{Concept} & \textbf{Tokens} & \textbf{Price / 1M} & \textbf{Cost} \\
					\midrule
					Input  & \num{114191269} & \$0.03 & \$3.43 \\
					Output & \num{30409898}  & \$0.13 & \$3.95 \\
					\midrule
					\textbf{Total} & \textbf{\num{144601167}} & & \textbf{\color{accent}\$7.38 USD} \\
					\bottomrule
			\end{tabular}}
		\end{center}
		\mynote{At API list pricing, the 200{,}000-call run would have cost \$7.38.
			Our self-hosted vLLM setup (RunPod GPUs), at \$1.04/hour, ran the
			same workload in 140--150 minutes for \$2.43--\$2.60.}
	\end{frame}
	
	% ================================================================
	% SLIDE -- Prompt anatomy (input / output example)
	% ================================================================
	\begin{frame}{Prompt anatomy: input and output}
		\begin{columns}[T]
			\column{0.5\textwidth}
			\textbf{Input} \; \footnotesize(persona + election context)
			\vspace{0.2em}
			\begin{block}{}
				\scriptsize
				\texttt{DEMOGRAPHIC PROFILE:}\\
				\texttt{Sex: Female \quad Age: 65}\\
				\texttt{State: North Carolina}\\
				\texttt{Race/Ethnicity: Black \quad \dots}\\
				\texttt{Education / Occupation: \dots}\\[0.4em]
				\texttt{ELECTION CONTEXT: \dots}\\
				\texttt{TASK: estimate voting preference}\\
				\texttt{FORMAT: valid JSON, sums to 100}
			\end{block}
			
			\column{0.5\textwidth}
			\textbf{Output} \; \footnotesize(JSON with reasoning)
			\vspace{0.2em}
			\begin{block}{}
				\scriptsize
				\texttt{\{}\\
				\texttt{\ "Donald Trump": 15,}\\
				\texttt{\ "Kamala Harris": 85,}\\
				\texttt{\ "reason": "Elderly Black woman}\\
				\texttt{\ from NC leaning Democratic on}\\
				\texttt{\ social issues favors Harris.",}\\
				\texttt{\ "in\_tok": 483, "out\_tok": 97}\\
				\texttt{\}}
			\end{block}
		\end{columns}
	\end{frame}
	
	% ================================================================
	% EXPERIMENT I -- V1: demographics only  (images: demo1, demo2)
	% ================================================================
	\begin{frame}{Experiment 1: Demographic}
		Baseline prompt uses \textbf{only} the demographic profile:
		sex, age, state, race/ethnicity, education, occupation.
	\end{frame}
	
	\begin{frame}{Experiment 1: Demographic --- predicted vs.\ real (2024)}
		\bigimg{demo1.png}
	\end{frame}
	
	\begin{frame}{Experiment 1: Demographic --- classification accuracy}
		\bigimg{demo2.png}
	\end{frame}
	
	% ================================================================
	% EXPERIMENT II -- V2: + cultural background  (images: cu1..cu4)
	% ================================================================
	\begin{frame}{Experiment 2: Cultural}
		Same skeleton as V1, adding a \textbf{cultural background} block.
		\begin{center}
			\color{amber}\itshape
			The cultural background prompt shifts the simulated electorate
			toward \textbf{more Republican votes} than the demographic-only
			baseline.
		\end{center}
	\end{frame}
	
	\begin{frame}{Experiment 2: Cultural --- predicted vs.\ real (2024)}
		\bigimg{cu1.png}
	\end{frame}
	
	\begin{frame}{Experiment 2: Cultural --- classification accuracy}
		\bigimg{cu2.png}
	\end{frame}
	
	\begin{frame}{Experiment 2: Cultural --- vote share and electoral vote changes}
		\splitimg{cu3.png}{cu4.png}
		\vspace{0.5em}
		\begin{center}
			\color{amber}\itshape
			Trump wins \textbf{more electoral votes} in the cultural-background
			prompt than in the demographic-only baseline.
		\end{center}
	\end{frame}
	
	% ================================================================
	% EXPERIMENT III -- V3: + persona overview  (images: pe1..pe4)
	% ================================================================
	\begin{frame}{Experiment 3: Persona}
		Same skeleton as V2, adding a full \textbf{persona overview}.
	\end{frame}
	
	\begin{frame}{Experiment 3: Persona --- predicted vs.\ real (2024)}
		\bigimg{pe1.png}
	\end{frame}
	
	\begin{frame}{Experiment 3: Persona --- classification accuracy}
		\bigimg{pe2.png}
	\end{frame}
	
	\begin{frame}{Experiment 3: Persona --- vote share and electoral vote changes}
		\splitimg{pe3.png}{pe4.png}
		\vspace{0.5em}
		\begin{center}
			\color{amber}\itshape
			The full persona overview produces the \textbf{largest shift toward
				Harris} of all four experiments.
		\end{center}
	\end{frame}
	
	% ================================================================
	% EXPERIMENT IV -- V4+V5: career goals + Big Five (OCEAN)  (images: ca1..ca4)
	% ================================================================
	\begin{frame}{Experiment 4: Career + OCEAN --- predicted vs.\ real (2024)}
		\bigimg{ca1.png}
		\begin{center}
			\color{amber}\itshape
			Adding career goals and OCEAN personality essentially
			\textbf{ties} the two candidates.
		\end{center}
	\end{frame}
	
	\begin{frame}{Experiment 4: Career + OCEAN --- classification accuracy}
		\bigimg{ca2.png}
	\end{frame}
	
	\begin{frame}{Experiment 4: Career + OCEAN --- vote share and electoral vote changes}
		\splitimg{ca3.png}{ca4.png}
	\end{frame}
	
	% ================================================================
	% EL SALVADOR
	% ================================================================
	\begin{frame}{Replication: El Salvador}
		\textbf{Setup:} proportional-to-population allocation, \textbf{100{,}000} or \textbf{50{,}000}
		synthetic personas, same 3-prompt pipeline as the U.S.\ core
		experiments (Demographic, Cultural, Persona) Datasets only have \textbf{148{,}000} .
		\vspace{0.4em}
		
		\textbf{Also tested:}
		\begin{itemize}
			\item \textbf{Language:} Spanish vs.\ English demographic prompt.
			\item \textbf{Topic-informed prompt:} top issues per survey\footnotemark:
			\emph{economy}, \emph{violence}, \emph{governance}.
		\end{itemize}
		\footnotetext{\tiny FUNDAUNGO/CEOP, \textit{Panel Electoral: El Salvador
				2022--2024} (2024).}
	\end{frame}
	
	% ================================================================
	% BRAZIL 2026
	% ================================================================
	\begin{frame}{Next test: Brazil 2026 presidential election}
		First round \textbf{4 Oct} $\cdot$ second round \textbf{25 Oct 2026}.
		Same 3-prompt setup as El Salvador (Demographic, Cultural, Persona).
		\vspace{0.4em}
		
		\textbf{Also tested:} Portuguese vs.\ English demographic
		prompt; topic-informed prompt using top issues per survey\footnotemark:
		\emph{violence}, \emph{economy}, \emph{corruption}, \emph{health}.
		\vspace{0.4em}
		
		\textbf{Second round:} current polling is \textbf{tight}:
		\textbf{Lula 43\%} vs.\ \textbf{Flávio Bolsonaro 40\%}. We generate an runoff prediction.
		\footnotetext{\tiny AS/COA, \textit{Poll Tracker: Brazil's 2026
				Presidential Election} (2026).}
	\end{frame}
	
\end{document}